\chapter*{How to Use This Book}
\addcontentsline{toc}{chapter}{How to Use This Book}
\markboth{How to Use This Book}{How to Use This Book}

\section*{What kind of book is this?}

This is neither an encyclopedic textbook nor a private formula sheet. It is a structured set of undergraduate mathematics notes designed for two uses:
\begin{enumerate}
    \item a \textbf{guided review}, when the reader has seen a subject before but wants clearer connections, proofs, and examples;
    \item a \textbf{reference map}, when the reader needs to locate a definition, theorem, technique, or application and understand how it fits into a larger theory.
\end{enumerate}
A complete first course still benefits from lectures, problem classes, and a specialized textbook. The exercises in this manuscript are intended to make reading active, not to imitate an entire semester of assessed work.

\section*{Recommended routes}

\begin{description}[leftmargin=3.3cm,style=nextline]
    \item[Standard undergraduate route] Begin with Chapter 2, Mathematical Analysis I, and Chapter 3, Linear Algebra. Continue to Chapters 4--7 according to interest. Treat Chapters 0--1 as optional reference material.
    \item[Foundations route] Read Chapters 0--1 before Chapter 2. This route is appropriate when the construction of number systems, formal logic, and axiomatic set theory are themselves part of the objective.
    \item[Computation and data route] Read Chapter 3, then Chapters 8 and 9. Return to Chapters 2 and 5 for optimization, approximation, integration, and Fourier analysis.
    \item[Structure route] Read Chapters 3, 4, 6, and 7. This route emphasizes vector spaces, algebraic structures, topological invariants, and holomorphic structure.
\end{description}

\section*{Dependency map}

Most chapters are locally self-contained, but the following dependencies are useful:
\[
\begin{array}{ccl}
\text{Analysis II} &\longleftarrow& \text{Analysis I},\\
\text{Abstract Algebra} &\longleftarrow& \text{basic Linear Algebra},\\
\text{Topology} &\longleftarrow& \text{sets, functions, and elementary analysis},\\
\text{Complex Analysis} &\longleftarrow& \text{Analysis I and basic topology},\\
\text{Probability} &\longleftarrow& \text{calculus; measure theory sections are optional on first reading}.
\end{array}
\]
Graph Theory can be read almost independently.

\section*{Proof and scope conventions}

A theorem followed by a \textbf{proof} is intended to contain a complete argument at the level of the chapter. A result described as an \textbf{advanced theorem} is included for orientation when a full proof would require substantial additional machinery. No omitted argument should be mistaken for an elementary exercise.

Definitions are followed, whenever useful, by examples, non-examples, or a computational interpretation. Application boxes show how a concept enters a model in science, engineering, or data analysis. They are not claims that the application is exhausted by the mathematics shown there.

\section*{How to work the exercises}

The exercises are arranged roughly from direct verification to synthesis. A good minimum routine is:
\begin{enumerate}
    \item reproduce one proof from memory;
    \item solve two computational exercises;
    \item solve one conceptual or counterexample exercise;
    \item explain one application in your own words.
\end{enumerate}
Selected hints and answers are intentionally brief. They indicate a route or a check, not a substitute for a written solution.

\begin{mynote}
The foundational material is available when it clarifies a later construction, but it should not become an entrance fee for learning calculus, linear algebra, graph theory, or probability. Start where your mathematical question begins, and move backward only when a dependency genuinely matters.
\end{mynote}
